\section{Introduction}

Serializable snapshot isolation~\cite{ports2012ssi} achieves serializability
without read locks by detecting, at commit time, the dependency patterns that
could produce a non-serializable schedule. One of the dependencies it must see is
between a transaction that read a range of values and a transaction that inserted
a value into that range~--- the phantom. The reader did not see the new row, and
nothing in its read set records that it \emph{would} have seen it, because the
row did not exist.

The mechanism is a predicate lock: the reader registers the region it examined,
the writer checks whether its new value falls inside any registered region. What
makes this cheap or expensive is granularity. An access method that offers no
support forces the region to be the entire relation, and then every insertion
conflicts with every scan, which turns serializable isolation from a correctness
guarantee into an availability problem.

For a B-tree the granularity question has an obvious answer, because a scan reads
a contiguous interval of a total order and the interval is exactly the predicate.
For an extensible index framework there is no order, the key type is not known to
the core, and what a scan examined has no description outside the structure
itself. This paper is about what granularity is available in that setting, and
about a case where the natural answer for one structure was carried over to
another where its justification did not hold.

\paragraph{Contributions.}

\begin{enumerate}
\item Page-granularity predicate locking for the generalized search tree, with
      the covering argument that establishes its correctness
      (Section~\ref{sec:gist}).
\item An account of two defects in the corresponding mechanism for the
      generalized inverted index, and the revised principle
      (Section~\ref{sec:gin}).
\item The observation that a concurrency-control decision must be recorded in
      the description of the structure and not only in code, with the reason
      (Section~\ref{sec:writing-it-down}).
\end{enumerate}

\section{Page granularity for the generalized search tree}
\label{sec:gist}

A scan of a generalized search tree visits a set of pages, and on each page it
evaluates \code{consistent} against every entry. What did it read? Not the
entries it returned, and not the entries it examined: it read the \emph{region}
described by the keys of the pages it visited, because had a matching entry been
present anywhere in that region, the scan would have found it.

That observation gives the granularity. A scan takes a predicate lock on every
page it examines; an insertion checks for a lock on the page where its entry will
be placed. The rule is correct because of a property the structure maintains
anyway: the key of a page covers every entry on the page. Hence if a new value
would be placed on a given page, it falls within that page's key, that is within
a region the scan of that page examined. There is no way for an insertion to
land on a page and yet be outside what a scan of that page would have covered.

Two implementation subtleties are worth recording because both are easy to get
wrong in the direction of false conflicts or missed ones. The check must not fire
during index construction, where entries are placed without any concurrent reader
having a stake in them. And splitting a node inserts a downlink into the parent,
which is an insertion into a page but not an insertion of user data; treating it
as the latter produces conflicts that correspond to nothing in the user's
schedule.

Implemented in \code{3ad55863e93}, released in PostgreSQL~11.

\section{Where the argument does not carry over}
\label{sec:gin}

The inverted index stores, for each key, the list of rows containing it: an entry
tree over keys, whose leaves hold either an inline posting list or the root of a
posting tree. Page granularity was implemented for it by analogy. The analogy is
where the trouble started, because the covering argument of
Section~\ref{sec:gist} does not hold: a page of the entry tree is not a region of
a space that an insertion must fall into, and a scan of this structure does not
examine a contiguous region at all.

Two defects followed, and they have opposite shapes.

\paragraph{A scan that finds nothing must still lock something.} A scan that
found no matching entry recorded no predicate lock. This is the natural reading
of ``lock what you read'' and it is wrong, because the purpose of a predicate
lock is to lock the \emph{gap}. A query asking for a key that does not exist has
read exactly the fact that it does not exist, and an insertion creating it must
conflict. Recording nothing means the phantom is invisible.

\paragraph{Turning on deferred insertion mid-flight bypasses existing locks.} The
index supports a pending list into which insertions may be deposited and applied
later. If that mode is enabled while transactions are running, subsequent
insertions go to the pending list and no longer conflict with locks that were
already taken on entry tree pages~--- the locks describe a structure the
insertion no longer touches. A configuration change thus silently weakens
isolation for transactions already in progress.

\paragraph{The revised principle.} A scan always locks the metapage, and this
lock represents the reading of the pending list whether or not one currently
exists~--- which is what makes the mode change harmless. And a scan always locks
the entry tree page it consulted, including when it found no match, which is what
makes gaps visible. In the other direction, locking was found to be excessive
and was reduced: every element of a posting tree belongs to a single key value,
so locking the root of that tree is sufficient and locking its pages
individually gains nothing.

Implemented in \code{0bef1c0678d}, released in PostgreSQL~11. A cost of the
mechanism showed up later in an unexpected place: the additional locking slowed
the isolation test suite enough under instrumentation to be worth
reporting~\cite{thread-predicate-gin-valgrind}, which is a reminder that
correctness machinery is paid for on every scan, not only on the schedules it
exists to catch.

\section{Writing the principle down where it can be checked}
\label{sec:writing-it-down}

The revision above was recorded in the prose description of the structure that
accompanies the source, not only in the code implementing it. We consider this
part of the result rather than a matter of documentation hygiene, for a reason
specific to concurrency control.

A rule such as ``a scan always locks the metapage, and this represents reading
the pending list whether or not one exists'' is not visible in the code that
implements it. What the code shows is a lock being taken at a particular line;
why it is taken there, and what would break if it were taken elsewhere or made
conditional, is nowhere in the text. A later change that makes the lock
conditional~--- for instance because profiling shows the pending list is usually
absent~--- looks locally like a valid optimisation, passes every functional test,
and reintroduces exactly the defect described above. The decision has to be
written where the person making that later change will encounter it.

This generalizes past predicate locking, and it is the same argument we have made
elsewhere about verifying structural invariants: properties that no functional
test observes must be recorded in a form that can be checked by something other
than running the code. An invariant checker is one such form. A stated principle
in the description of the structure is a weaker one, since it is checked by a
human reader rather than by a program, but it is much cheaper and it is available
for properties that no program can practically check.

\section{Evaluation}
\label{sec:eval}

% TODO. План измерений; сборки парные, различающиеся ровно одним коммитом.
%
% 1. Ложные конфликты: доля транзакций, откатываемых из-за сериализационного
%    конфликта, при страничной гранулярности против блокировки отношения
%    целиком (3ad55863e93^ / 3ad55863e93). Нагрузка — чтение диапазона плюс
%    вставка в непересекающийся диапазон: при верной гранулярности конфликтов
%    быть не должно вовсе, при блокировке отношения — должны быть все.
%    Величина двоичная и потому малошумная.
% 2. Цена механизма: пропускная способность при SERIALIZABLE против
%    REPEATABLE READ на нагрузке без конфликтов, до и после изменения.
%    Это та величина, о которой сообщалось в треде о замедлении тестов.
% 3. Обратный индекс: воспроизведение обоих дефектов изоляционными тестами до
%    0bef1c0678d и их отсутствие после. Первичный результат здесь двоичный —
%    обнаруживается аномалия или нет, — а не числовой.

\section{Limitations}

Page granularity is coarser than the predicate a query actually describes, and
on a page holding entries spread over a wide region it will report conflicts that
a finer mechanism would not. We have not characterized how often that happens on
real workloads; the argument for page granularity is that it is sound and
cheaply computable from what the structure already maintains, not that it is
tight.

The correctness argument of Section~\ref{sec:gist} depends on the covering
property of keys, and therefore on the operator class maintaining it. An operator
class whose \code{union} does not produce a covering key would break predicate
locking along with search itself, so this is not an additional assumption; but it
is worth naming, since it is the point at which the framework trusts code it did
not write.

Section~\ref{sec:gin} reports defects found in production and by review rather
than a systematic search. We do not claim the revised principle is complete, and
the history of this particular mechanism does not encourage the claim.

\section{Conclusion}

Predicate locking is the point where an access method has to describe what a
query read, rather than what it returned, and the two differ exactly by the rows
that do not exist. For a structure with a covering invariant the description is
available for free: the key of a page is a sound over-approximation of what a
scan of that page examined, so page granularity is both correct and cheap.

The instructive part is the failure of the analogy. The same granularity was
applied to a structure without that invariant, and produced a mechanism that
looked equivalent and was not: it recorded nothing when a scan found nothing, and
it could be bypassed by a configuration change. Neither defect is visible in a
test that runs a workload and checks its answers, because both concern schedules
that the test does not generate. What they have in common with the rest of our
work on this framework is the form of the eventual remedy: state the property, in
a place where a later change will be checked against it, rather than trusting
that the code will keep saying what it currently says.
